\documentclass[aps,prl,twocolumn,superscriptaddress,longbibliography]{revtex4-2}
\usepackage{amsmath,amssymb,bm}
\usepackage{graphicx}
\usepackage{hyperref}
\usepackage{xcolor}

\begin{document}

\title{Phase-Sensitive Andreev Diagnostics for Unconventional Pairing in Twisted Bilayer WSe$_2$}

\author{Author A}
\affiliation{Institute Placeholder}
\author{Author B}
\affiliation{Institute Placeholder}
\author{Author C}
\affiliation{Institute Placeholder}

\date{May 1, 2026}

\begin{abstract}
Twisted bilayer WSe$_2$ has recently emerged as a non-graphene moir\'e
superconductor with a phase diagram shaped by strong correlations, displacement
field tuning, and twist-angle evolution. The pairing symmetry, however, remains
unresolved because ordinary tunnelling features can mix superconducting and
correlation-driven energy scales. We develop a generalized
phase-sensitive Andreev framework for twisted bilayer WSe$_2$ in which the
normal correlated background and the anomalous superconducting self-energy are
kept distinct from the outset. The resulting theory yields three main results.
First, Andreev response selectively tracks the inner superconducting scale,
whereas high-barrier tunnelling retains sensitivity to outer correlated
features. Second, ordinary same-sign $s$ wave fails structurally once barrier,
field angle, and interface orientation are constrained on the same
$\alpha$-$Z$-$\phi$ grid. Third, finite-bias peaks alone are not a sufficient
indicator of chirality: only the chiral branch exhibits phase-locked drift of
the peak position together with $\alpha$-coupled peak splitting, whereas a
phase-blind symmetric double-peak control reproduces finite-bias peaks but not
their phase-sensitive evolution. These results define an experimentally usable
decision tree for separating ordinary pairing, sign-changing channels, and a
chiral candidate family in twisted bilayer WSe$_2$.
\end{abstract}

\maketitle

Twisted moir\'e materials provide a route to superconductivity in flat bands
beyond graphene, but they also compress correlation scales and pairing scales
into the same narrow energy window~\cite{Cao2018SC,Cao2018CI,Yankowitz2019,Lu2019,Park2021,Balents2020}.
The problem is especially severe in semiconducting moir\'e transition-metal
dichalcogenides, where filling, bandwidth, valley structure, and interaction
content all shift under the same electrostatic controls~\cite{Tang2020,Regan2020,Xu2022BilayerHubbard}.

Twisted bilayer WSe$_2$ is now the sharpest test case in this class. Robust
superconductivity has been established near $3.5^\circ$-$3.65^\circ$ and again
near $5.0^\circ$, and later measurements connected these windows through angle
evolution and bandwidth-tuned Mottness within one materials
family~\cite{Xia2025SC,Hone2025SC5,Angle2026,Bandwidth2026}. Spectroscopy on
moir\'e superconductors has also made a crucial point: outer correlation-driven
features and inner superconducting gaps need not be probed by the same
transport channel~\cite{Kim2026Intervalley}. A pairing claim based only on the
presence of a low-energy peak is therefore not yet diagnostic.

The opportunity is therefore to use Andreev spectroscopy as a phase-sensitive
probe rather than as a refined fitting exercise. In conventional BTK theory,
subgap conductance is set by the competition among Andreev reflection, normal
reflection, and interface transparency~\cite{BTK1982,Octavio1983,Beenakker1997}.
For unconventional superconductors, the same interface problem can project the
order-parameter phase structure into bound-state formation, peak shifts, and
anisotropic conductance~\cite{Hu1994,TanakaKashiwaya1995,KashiwayaTanaka2000,Linder2010,Wu2018Shift,Wang2023STM}.
The relevant question for twisted bilayer WSe$_2$ is therefore not whether
finite-bias structure can be produced, but whether that structure must inherit
phase-sensitive locking among field angle, interface orientation, and barrier
evolution.

Here we formulate a generalized phase-sensitive Andreev framework for twisted
bilayer WSe$_2$ and use it to build a decision tree for pairing
identification. We first separate the correlated normal background from the
anomalous superconducting channel and show that Andreev response selectively
tracks the inner superconducting scale. We then compare ordinary same-sign
$s$ wave, an intervalley sign-changing channel, and a chiral candidate family
on the same $\alpha$-$Z$-$\phi$ grid. Finally, we introduce a non-chiral
control kernel with symmetric peaks at $\pm V$ to test whether finite-bias
peaks alone diagnose chirality. They do not: only the chiral branch shows
phase-locked peak-position drift together with $\alpha$-coupled peak splitting,
which the control channel cannot reproduce.

\begin{figure}[t]
  \centering
  % \includegraphics[width=\columnwidth]{fig1_schematic.pdf}
  \fbox{\parbox[c][4.8cm][c]{0.95\columnwidth}{\centering
  Figure 1 schematic placeholder:\\
  minimal normal-state model, separation of $\Sigma_{\rm corr}$ and $\Delta$,
  candidate pairing families, and measurement geometry for the
  $\alpha$-$Z$-$\phi$ scan.}}
  \caption{\textbf{Minimal model and observable hierarchy.}
  Schematic-only figure. The normal correlated background
  $\Sigma_{\rm corr}$ and the anomalous superconducting sector $\Delta$ are
  separated at the Hamiltonian level so that high-barrier tunnelling and
  Andreev-dominated conductance can be compared without assuming a single common
  energy scale. The figure also defines the three candidate pairing families
  used below---ordinary same-sign $s$ wave, an intervalley sign-changing
  channel, and a chiral family---together with the interface orientation
  $\alpha$, barrier strength $Z$, and in-plane field angle $\phi$ that enter the
  generalized BTK analysis.}
  \label{fig:model}
\end{figure}

The starting point is a low-energy moir\'e Hamiltonian in which the normal-state
sector, valley structure, and correlated background are retained separately from
the anomalous self-energy,
\begin{equation}
H_{\rm BdG}(\mathbf{k})=
\begin{pmatrix}
H_N(\mathbf{k})-\mu+\Sigma_{\rm corr}(E,\mathbf{k}) & \Delta(\mathbf{k}) \\
\Delta^\dagger(\mathbf{k}) &
-H_N^{T}(-\mathbf{k})+\mu-\Sigma_{\rm corr}^{T}(-E,-\mathbf{k})
\end{pmatrix}.
\end{equation}
This separation is not cosmetic. It is required if the tunnelling channel is to
retain sensitivity to outer correlated features while the Andreev channel is
allowed to track the anomalous sector more directly. Guided by recent
experimental and theoretical work on twisted WSe$_2$ and related moir\'e
platforms, we keep three candidate families in the reduced pairing library:
ordinary same-sign $s$ wave, an intervalley sign-changing channel, and a chiral
family representing a phase-winding unconventional branch~\cite{Schrade2024,TopologicalWSe22025,NuOne2025,ChiralBandsMoTe22025,ReadGreen2000,SigristUeda1991,SatoAndo2017}.

\begin{figure}[t]
  \centering
  % \includegraphics[width=\columnwidth]{fig2_selectivity.pdf}
  \fbox{\parbox[c][5.0cm][c]{0.95\columnwidth}{\centering
  Figure 2 data placeholder:\\
  high-barrier tunnelling versus Andreev-dominated conductance;\\
  $E_{\rm out}$, $\Delta_{\rm in}$, and $E_{\rm AR}$ across control parameters.}}
  \caption{\textbf{Inner-gap selectivity of the Andreev channel.}
  High-barrier tunnelling retains the outer correlated feature $E_{\rm out}$,
  whereas the Andreev-dominated response follows the inner superconducting scale
  $\Delta_{\rm in}$. The key statement is channel selectivity, not a two-peak
  fit: within the same normal-state benchmark, $E_{\rm AR}$ tracks the anomalous
  scale more closely than the outer correlated feature. This establishes the
  transport logic needed for the later phase-sensitive analysis.}
  \label{fig:selectivity}
\end{figure}

Figure~\ref{fig:selectivity} summarizes the first step of that logic. Once
$\Sigma_{\rm corr}$ and $\Delta$ are separated at the model level, the
high-barrier tunnelling response retains the outer correlated scale
$E_{\rm out}$, while the Andreev-dominated response follows the inner
superconducting scale $\Delta_{\rm in}$. This distinction is the minimal
evidence required before any pairing claim can be made. Without it, low-energy
structure in the conductance could still be reinterpreted as a correlated
normal-state feature rather than a superconducting diagnostic~\cite{Kim2026Intervalley}.

\begin{figure}[t]
  \centering
  % \includegraphics[width=\columnwidth]{fig3_fingerprints.pdf}
  \fbox{\parbox[c][5.0cm][c]{0.95\columnwidth}{\centering
  Figure 3 data placeholder:\\
  generalized BTK fingerprints on the $\alpha$-$Z$-$\phi$ grid and their
  reduction to a small set of observables.}}
  \caption{\textbf{Phase-sensitive fingerprints on the $\alpha$-$Z$-$\phi$
  grid.} Barrier evolution, interface orientation, and in-plane field angle are
  scanned on the same grid for the three candidate pairing families. Ordinary
  same-sign $s$ wave remains comparatively smooth, whereas unconventional
  channels develop stronger barrier sensitivity and orientation dependence. The
  figure should be read through reduced observables such as subgap average,
  finite-bias peak position, and peak splitting rather than through curve
  counting alone.}
  \label{fig:fingerprints}
\end{figure}

The generalized BTK scan on the common $\alpha$-$Z$-$\phi$ grid then separates
two distinct types of unconventional response. The sign-changing channel first
breaks the monotonic barrier trend that characterizes ordinary $s$ wave under
the same constraints, thereby establishing the exclusion step. The chiral branch
goes further by reorganizing the interface spectrum around finite-bias edge
states, so that the dominant spectral weight no longer remains pinned to
zero bias. Figure~\ref{fig:fingerprints} therefore serves as the bridge between
the scale-selectivity logic of Fig.~\ref{fig:selectivity} and the control-based
chirality test of Fig.~\ref{fig:control}.

\begin{figure}[t]
  \centering
  % \includegraphics[width=\columnwidth]{fig4_control_discriminant.pdf}
  \fbox{\parbox[c][5.2cm][c]{0.95\columnwidth}{\centering
  Figure 4 data placeholder:\\
  chiral branch versus phase-blind symmetric double-peak control;\\
  positive-peak drift, peak splitting, and barrier evolution.}}
  \caption{\textbf{Control discriminant for chiral finite-bias peaks.}
  Comparison between the chiral finite-bias edge-state kernel and a non-chiral
  control kernel with symmetric peaks at $\pm V$, tunable width, and tunable
  spectral weight, evaluated on the same $\alpha$-$Z$-$\phi$ grid. The control
  channel generates finite-bias double peaks but remains phase blind: the
  positive-bias peak position does not drift with $\phi$, and the peak splitting
  does not couple to $\alpha$. By contrast, the chiral channel exhibits both
  phase-locked peak-position drift and $\alpha$-coupled peak splitting. In the
  present scan, the maximum positive-peak drift is $0.0928$ for the chiral
  channel and $0$ for the control, while the maximum $\alpha$-coupled change in
  peak splitting is $0.1856$ for the chiral channel and $0$ for the control.
  The chiral branch is also non-monotonic over the full tested configuration
  set, whereas the control channel retains a monotonic fraction of $0.67$.
  Finite-bias peaks alone are therefore insufficient: the manuscript claim rests
  on the phase-sensitive coupling among peak position, interface orientation,
  and barrier evolution, which the trivial double-peak control fails to
  reproduce.}
  \label{fig:control}
\end{figure}

Figure~\ref{fig:control} supplies the control test that converts finite-bias
structure from a suggestive spectral feature into a phase-sensitive statement.
The non-chiral control kernel reproduces finite-bias double peaks, but it
remains phase blind: the positive-bias peak position shows no $\phi$-dependent
drift, and the peak splitting shows no coupling to $\alpha$. The chiral branch
shows both on the same grid. In the present scan, the maximum drift of the
positive-bias peak position reaches $0.0928$, while the maximum
$\alpha$-coupled change in peak splitting reaches $0.1856$; both quantities
vanish in the control channel. The chiral branch is also non-monotonic
throughout the tested configuration set, whereas the control channel retains a
monotonic fraction of $0.67$. Finite-bias peaks are therefore not enough by
themselves. What distinguishes the chiral branch is the phase-sensitive locking
among peak position, interface orientation, and barrier evolution.

Taken together, Figs.~\ref{fig:selectivity}--\ref{fig:control} support a
layered conclusion. The present evidence is strong enough to exclude ordinary
same-sign $s$ wave and to identify an unconventional family. The control test
then isolates what is special about the chiral branch within that family: not
the existence of finite-bias peaks, but their phase-sensitive coupling to
$\phi$, $\alpha$, and $Z$. This remains more conservative than a full
topological identification, which would require an additional invariant,
phase-bias, or independent boundary-mode benchmark~\cite{ReadGreen2000,FuBerg2010,SatoAndo2017}.

The remaining limitations are equally clear. The present draft still relies on
the strongest available compressed benchmark rather than on a fully rerun
true-Tuo-TB figure set, and the BTK kernel remains a physically motivated
reduced model rather than a complete multiorbital scattering solver. These
limitations do not erase the control discriminant developed here, but they do
set the boundary of the claim. What is established at this stage is a
phase-sensitive diagnostic structure that excludes ordinary pairing and
separates a chiral candidate from a trivial finite-bias control on the same
observable grid.

\bibliographystyle{apsrev4-2}
\bibliography{prl-manuscript-draft-2026-05-01}

\end{document}
